\section{封国、关中与恢复：战争结束以后}
\label{sec:han-founding-recovery}

前二〇六年的关中，并不是一个等待新皇帝接收的空壳。秦朝已亡，仓储、关隘、县廷和道路却仍在；项羽以十八王重排这些资源，刘邦被徙至汉中。\eventref{WH-0206-EIGHTEEN-KINGS}{十八王分封}表面上兑现军功，实际把战后最重要的地区拆给不同的军政集团。封号能给各方一个名义，却不能规定粮食从哪里来、旧吏听谁的命令、军队是否愿意久驻异地。分封没有封住竞争，反而把它铺在秦的旧行政网络上。

刘邦出汉中、还定三秦，获得的首先不是一个抽象的“根据地”，而是关中能供给军队和容纳官署的条件。彭城败后，荥阳—成皋的狭长通道、河北赵地的兵员、齐地的粮赋和关中后方被拉进同一场战争。垓下的合围因而不是某一位统帅凭一计制胜，而是多个战区终于能够同时向西楚施压。\eventref{WH-0202-HAN-FOUNDATION}{汉朝建立}把军事联盟置于皇帝名号之下；它尚未解决联盟成员各自拥有的兵众、食邑和地方关系。

\subsection{长安不是终点，而是调度的起点}

定都长安，意味着新朝把关中作为文书、粮食和军队汇聚的中枢。长安靠近秦旧的仓储与险要，也远离关东的功臣网络；这项选择同时给了皇帝一个后方，也提醒朝廷必须向东部持续输送诏令、赏赐和任命。\eventref{WH-0200-BAIDENG}{白登之围}进一步揭出这个中枢的限度：北方骑兵、道路和补给并非关中一声征发便能取得。和亲、互市与边郡防御不是软弱或强硬的简单标签，而是在骑兵条件和财力不足时保存选择的组合。

异姓王的处置也应放在这个压力下理解。韩信、彭越、英布等人曾为战胜西楚提供兵力和地方联系；一旦皇帝需要稳定征发，他们又成了能够自行集结人马的风险。徙封、废杀和以刘氏宗室替代异姓王，确实让皇帝更容易将军政资源接回中枢，却不等于地方从此失去自己的门户。早汉形成的不是纯粹郡县国家，而是一张由皇帝、侯王、郡县吏员、功臣家族和边郡将领共同承重的网。

\subsection{恢复并不只写在诏令上}

吕后临朝与前一八〇年的诛吕，说明继承本身就是重新分配这张网的时刻。功臣、宗室、宫廷军力和外戚并未按照血缘自动站队；他们要判断谁能控制禁军、谁能给出可接受的名分、谁能保证既有的封爵和任命。文帝朝减田租、弛山泽禁、调整刑罚，常被后来的叙事统称为休养生息。诏令当然重要，但它首先告诉我们朝廷希望怎样恢复可登记的土地与人户；它不能逐户报出劳役、转输、欠赋和地主关系的实际重量。

放宽民间铸钱尤其能看出这一点。田租面向土地与编户，钱币则把矿产、范铸技术、商人和王国财力带进另一种流通。政策松动并不是国家消失，而是国家在资源不足时改变了直接经营和间接控制的比例。到景帝时，宗室王国仍有封号、官属和地方声望，中央又已无法容忍它们保持可独立调兵的空间。\eventref{WH-0154-CUOFAN}{七国之乱}的爆发，将这项长期矛盾从奏议与削地程序推到道路、粮道和军队的争夺上。

周亚夫断吴楚粮道的胜利，并没有把王国从地图上抹去。它改变的是中央核验王国任官、司法与军政的能力。新朝从战后恢复中取得的真正资源，不只是一批更低的税率或更安静的年份，而是让名号、户籍、运输与军队开始彼此接合的若干接口；这些接口日后既可扩张，也会被外戚、边将和地方势力重新占用。

\begin{sourcenote}[从战场叙事到恢复史]
《史记》的本纪与列传擅长保存人物和战场的动作，《汉书》更便于核对建国后的年序、封国与制度位置。二者都不能单独把鸿门会的言辞、垓下的楚歌、白登的兵数或诛吕的密谋复原为现场。关中城址、钱范与后出的编年材料可补充不同尺度，但也各有保存范围。本节使用它们追踪资源如何被接合，而不把后来的“文景之治”倒写为开国时已经实现的社会状态。
\end{sourcenote}
